\documentclass[11pt]{article}

\usepackage[a4paper,margin=29mm]{geometry}
\usepackage{amsmath,amssymb,amsthm,mathtools}
\usepackage{enumitem}
\usepackage{microtype}
\usepackage{xcolor}
\usepackage[colorlinks=true,linkcolor=blue!55!black,citecolor=blue!55!black,
  urlcolor=blue!55!black]{hyperref}

\newtheorem{theorem}{Theorem}[section]
\newtheorem{proposition}[theorem]{Proposition}
\newtheorem{lemma}[theorem]{Lemma}
\newtheorem{corollary}[theorem]{Corollary}
\theoremstyle{definition}
\newtheorem{definition}[theorem]{Definition}
\theoremstyle{remark}
\newtheorem{remark}[theorem]{Remark}

\newcommand{\C}{\mathbb C}
\newcommand{\R}{\mathbb R}
\newcommand{\Z}{\mathbb Z}
\newcommand{\N}{\mathbb N}
\newcommand{\A}{\mathcal A}
\newcommand{\B}{\mathcal B}
\newcommand{\D}{\mathcal D}
\newcommand{\Hh}{\mathcal H}
\newcommand{\Q}{\mathcal Q}
\newcommand{\cK}{\mathcal K}
\newcommand{\cM}{\mathcal M}
\newcommand{\HC}{\operatorname{HC}}
\newcommand{\Kzero}{K_0}
\newcommand{\Str}{\operatorname{Str}}
\newcommand{\Tr}{\operatorname{Tr}}
\newcommand{\ev}{\operatorname{ev}}
\newcommand{\id}{\operatorname{id}}
\newcommand{\dd}{\,\mathrm d}
\newcommand{\ol}[1]{\overline{#1}}

\title{\textbf{Real corona classes and cyclic transgression}\\
\large A quaternionic Hardy globalization of the Riemann spectral module}
\author{Michael Haufschild}
\date{}

\begin{document}
\maketitle

\begin{abstract}
Let $B$ be the Blaschke inner factor of Riemann's completed zeta function in the
critical strip.  We prove that the upper and lower zero factors $B_+$ and $B_-$
form a corona pair.  This produces a canonical idempotent
\[
 e\in H^\infty(0<\Re s<1)/BH^\infty,
 \qquad e^2=e,
\]
which equals one on upper zeros and zero on lower zeros.  Complex conjugation
exchanges $e$ and $1-e$.  Consequently the multiplier $m=2e-1$ converts the
ordinary Real structure on the Hardy quotient $H^2/BH^2$ into a bounded
quaternionic structure $J=m\#$ satisfying $J^2=-1$.  This construction is
unconditional and does not use the Riemann hypothesis, simplicity of zeros, or
zero separation.

The idempotent defines an anti-Real class
$\kappa=[e]-[1-e]\in K_0$, of infinite order, and every absolutely summable
zero weight gives a continuous cyclic $0$-cocycle whose Chern--Connes pairing
with $\kappa$ is the oriented upper-minus-lower zero character.  For a completed
$L$-function with real zeros in the open strip, a third corona idempotent
$e_{\R}$ is fixed by the Real involution.  It is the exact obstruction to a
quaternionic structure on a cyclic scalar fibre, while the surviving partial
operator satisfies $J_0^2=-(1-e_{\R})$.  Its cyclic pairing records precisely
the real-zero contribution.

The functional equation gives a second involution.  In the crossed product by
this involution, $e$ and $1-e$ become equivalent and
$m=[eu,ue]$.  Thus the oriented class disappears simultaneously from ordinary
$K_0$ and ordinary cyclic degree zero.  It survives in graded cyclic homology:
$m$ anticommutes with the crossed-product generator, and a canonical
two-dimensional supertrace pairs with it as $2$.  Finally, a Mellin--Hardy
density theorem identifies $H^2/BH^2$ as the natural Hilbert globalization of
the scalar Connes--Consani--Marcolli/Meyer zeta quotient, and matches the cyclic
pairing with its finite-packet spectral characters.  The algebraic core has an
unconditional Lean~4 formalization.
\end{abstract}

\medskip
\noindent\textbf{Keywords.}
Noncommutative geometry; cyclic homology; Hardy space; corona theorem;
Riemann zeta function; Real $K$-theory; crossed product; spectral realization.

\section{Introduction}

Connes' spectral interpretation of the zeros of zeta replaces a hypothetical
self-adjoint operator by a cohomological quotient on the adeles class space
\cite{Connes1999}.  Meyer constructed a closely related virtual representation
whose distributional character is Weil's explicit formula \cite{Meyer2005}.
Connes, Consani, and Marcolli subsequently expressed the formula as a Lefschetz
trace formula on cyclic homology and isolated positivity of the trace pairing as
the characteristic-zero analogue of the last step in Weil's proof
\cite{CCM2007}.  Recent zeta-cycle and zeta-spectral-triple constructions retain
this quotient-and-trace viewpoint \cite{ConnesConsani2023,CCM2025}.

There is a persistent Real-structure problem inside this program.  The divisor
of a self-dual completed $L$-function carries two involutions: conjugation
$s\mapsto\bar s$ and the functional equation $s\mapsto1-s$.  Paugam obtained a
symplectic spectral pairing from these symmetries \cite{Paugam2009}, but the
upper/lower spectral split is not supplied there by a canonical bounded
projection.  Without such a projection, the formal sign change between the two
halves does not define a bounded quaternionic operator on a Hilbert
globalization, nor does it define a $K$-class that can be paired with cyclic
cohomology.

The analytic observation of this paper is that the arithmetic divisor has much
more separation after the exponential uniformization of the strip than its
Euclidean geometry suggests.  Put
\[
        \mathbb S=\{s\in\C:0<\Re s<1\},\qquad
        w=\exp(\pi i s).
\]
The map $s\mapsto w$ sends $\mathbb S$ to the upper half-plane.  Upper zeros go
exponentially to $0$ and lower zeros go exponentially to infinity.  Polynomial
zero counting therefore becomes an absolutely summable two-ended divisor.
This is enough to prove a uniform corona estimate for the two Blaschke
products.  Carleson's theorem then produces the missing bounded projection.

The resulting structure has three layers.

\begin{enumerate}[label=\textup{(\roman*)}]
\item On the commutative Hardy multiplier quotient, the corona idempotent gives
      an infinite-order anti-Real $K_0$ class and a quaternionic operator.
\item Real spectral fixed points produce a fixed $K_0$ defect and are exactly
      the fibres on which quaternionification fails.
\item Adjoining the functional equation makes the algebra noncommutative.  The
      ordinary orientation becomes a commutator, but a graded cyclic class
      remains and has an explicit nonzero supertrace pairing.
\end{enumerate}

This last transition is the cyclic transgression in the title.  It gives an
exact answer to a natural information question: which part of the oriented
zero character is lost when the functional equation is internalized as a
crossed-product symmetry, and where does that information survive?

\subsection{Main results}

We first formulate a divisor theorem that applies beyond zeta.  An admissible
Real divisor in $\mathbb S$ is a locally finite multiset $\Sigma$, invariant
under $s\mapsto\bar s$ and $s\mapsto1-s$, whose counting function is polynomially
bounded and which is the zero divisor of a function holomorphic across the
closed finite strip.  Its nonreal parts are denoted $\Sigma_+$ and $\Sigma_-$.
The last condition implies a positive distance between the nonreal divisor and
the real interval whenever the real part has been removed.

\begin{theorem}[Real-corona theorem, informal form]\label{thm:intro-corona}
Let $B=B_+B_-$ be the Blaschke product of a real-zero-free admissible Real
divisor.  Then
\[
          \inf_{s\in\mathbb S}\bigl(|B_+(s)|+|B_-(s)|\bigr)>0.
\]
There is a canonical idempotent $e$ in
$\A_B=H^\infty(\mathbb S)/BH^\infty(\mathbb S)$ with values $1$ on
$\Sigma_+$ and $0$ on $\Sigma_-$.  Both divisor involutions exchange $e$ and
$1-e$.
\end{theorem}

The word ``canonical'' refers to the quotient class.  Corona solutions are not
unique, but every solution gives the same element $e$ by the Chinese remainder
isomorphism.  No enumeration, interpolation basis, or zero-separation
hypothesis enters.

Let $\#$ be coefficient conjugation,
$f^\#(s)=\overline{f(\bar s)}$, and let it act on the Hardy quotient
$\Q_B=H^2(\mathbb S)/BH^2(\mathbb S)$.  The quotient has a canonical equivalent
Hilbert norm obtained from its Chinese remainder decomposition.

\begin{theorem}[Quaternionification and cyclic orientation, informal form]
\label{thm:intro-quaternion}
With $m=2e-1$, the anti-linear operator $J=m\#$ on $\Q_B$ is antiunitary for
the Chinese remainder norm and satisfies $J^2=-1$.  The class
\[
              \kappa=[e]-[1-e]\in K_0(\A_B)
\]
has infinite order and obeys $\kappa^\#=-\kappa$.  If a weight $h$ is absolutely
summable on $\Sigma$, then
\[
 \tau_h([f])=\sum_{\rho\in\Sigma}m_\rho h(\rho)f(\rho)
\]
is a continuous cyclic $0$-cocycle and
\[
 \langle[\tau_h],\kappa\rangle
 =\sum_{\rho\in\Sigma_+}m_\rho h(\rho)
  -\sum_{\rho\in\Sigma_-}m_\rho h(\rho).
\]
\end{theorem}

If a finite real-zero factor $B_{\R}$ is present, pairwise corona separation
produces $e_+,e_{\R},e_-$ with sum one.  Then $m_0=e_+-e_-$ satisfies
$m_0^2=1-e_{\R}$ and $J_0=m_0\#$ satisfies
\[
                     J_0^2=-(1-e_{\R}).                 \tag{1.1}
\]
The fixed class $\delta=[e_{\R}]$ is nonzero exactly when a real zero exists.
On the one-dimensional scalar fibre at such a zero no conjugate-linear map can
square to $-1$.  Thus $\delta$ is an exact obstruction, not merely a record of
where the construction was omitted.

Let $\sigma f(s)=f(1-s)$ and
$\B_B=\A_B\rtimes_\sigma\Z/2$, with generator $u$.

\begin{theorem}[Functional-equation transgression, informal form]
\label{thm:intro-crossed}
In $\B_B$, the elements $x=eu$ and $y=ue$ satisfy
\[
       xy=e,\qquad yx=1-e,
\]
so $\kappa$ maps to zero in $K_0(\B_B)$.  Moreover
\[
       m=xy-yx=[x,y],\qquad mu=-um.                    \tag{1.2}
\]
Every ordinary cyclic $0$-cocycle on $\B_B$ annihilates $m$.  For each
two-point functional-equation orbit there is, however, a graded representation
and a supertrace $\Str$ such that $\Str(m)=2$.  Hence the class of $m$ is zero in
ordinary $HC_0$ and nonzero in graded $HC_0$.
\end{theorem}

For Riemann's $\xi$-function the real defect is absent unconditionally.  We show
that its inner factor has no singular part, so Theorems
\ref{thm:intro-corona}--\ref{thm:intro-crossed} apply to the actual complete
nontrivial divisor.  A Mellin density theorem then connects the quotient to the
scalar CCM/Meyer source rather than introducing an unrelated model space.

\subsection{Scope and nonclaims}

The results do not imply the Riemann hypothesis.  Upper and lower refer to the
sign of $\Im\rho$, not to the side of the critical line.  The construction does
not produce a positive Weil form.  It also does not turn Meyer's nuclear
distributional trace into a trace-class Hilbert operator.  Schatten membership
for truncated Toeplitz operators requires additional Besov-type hypotheses
\cite{LopattoRochberg2016}; none is asserted here.  What is proved is a bounded
Real/quaternionic splitting, its $K$-theoretic and cyclic character, the exact
fixed-point obstruction, and the ordinary-to-graded change forced by the
functional equation.

\section{Hardy spaces and admissible divisors}

We use standard Hardy-space facts in the strip and upper half-plane; see
\cite{Garnett2007,Nikolski2002}.  Equip $H^2(\mathbb S)$ with the equivalent
boundary norm
\[
 \|F\|_{H^2_\partial}^2
 =\|F(0+i\,\cdot)\|_{L^2(\R)}^2
  +\|F(1+i\,\cdot)\|_{L^2(\R)}^2.                    \tag{2.1}
\]
Boundary values are understood non-tangentially.  Multiplication by
$H^\infty(\mathbb S)$ acts boundedly on this space.

The exponential map
\[
        \Phi(s)=e^{\pi i s}                            \tag{2.2}
\]
is a conformal bijection from $\mathbb S$ to the upper half-plane $\C_+$.
For $a\in\C_+$ write
\[
        b_a(w)=\frac{w-a}{w-\bar a}.                   \tag{2.3}
\]
Its boundary modulus is one and
\[
 1-|b_a(w)|^2
 =\frac{4\,\Im w\,\Im a}{|w-\bar a|^2}.              \tag{2.4}
\]

\begin{definition}
A multiset $\Sigma\subset\mathbb S$ is \emph{strip-admissible} if:
\begin{enumerate}[label=\textup{(A\arabic*)}]
\item it is the divisor in $\mathbb S$ of a nonzero function holomorphic on a
      neighbourhood of the closed finite strip;
\item $N_\Sigma(T)=\#\{\rho\in\Sigma:|\Im\rho|\le T\}=O((1+T)^d)$ for some
      $d$;
\item it is invariant, with multiplicities, under $\rho\mapsto\bar\rho$ and
      $\rho\mapsto1-\rho$.
\end{enumerate}
It is \emph{real-zero-free} if $\Sigma\cap\R=\varnothing$.
\end{definition}

Condition (A1) is deliberately stronger than local finiteness in the open
strip.  It excludes a hidden sequence accumulating at a finite boundary point.
Together with real-zero-freeness it implies
\[
       \inf_{\rho\in\Sigma}|\Im\rho|>0.               \tag{2.5}
\]

\begin{lemma}[Exponential two-ended Blaschke estimate]
\label{lem:bl-sum}
Let $\Sigma$ be strip-admissible and real-zero-free.  If
$a_\rho=e^{\pi i\rho}$, then for some $0<r<1$,
\[
 \begin{aligned}
 &|a_\rho|\le r\quad(\rho\in\Sigma_+),
 &&\sum_{\rho\in\Sigma_+}m_\rho|a_\rho|<\infty,\\
 &|a_\rho|\ge r^{-1}\quad(\rho\in\Sigma_-),
 &&\sum_{\rho\in\Sigma_-}m_\rho|a_\rho|^{-1}<\infty.
 \end{aligned}                                         \tag{2.6}
\]
In particular both halves define Blaschke products in $\C_+$.
\end{lemma}

\begin{proof}
Equation (2.5) gives the uniform inequalities with
$r=e^{-\pi\inf|\Im\rho|}$.  Divide the positive imaginary axis into intervals
$[n,n+1)$.  Polynomial zero counting bounds the number of zeros in the $n$th
interval by $O((n+1)^d)$.  Hence
\[
 \sum_{\rho\in\Sigma_+}e^{-\pi\Im\rho}
 \le C\sum_{n\ge0}(n+1)^d e^{-\pi n}<\infty.
\]
The lower-half assertion is identical because
$|a_\rho|^{-1}=e^{\pi\Im\rho}=e^{-\pi|\Im\rho|}$.
The usual upper-half-plane Blaschke condition follows a fortiori.
\end{proof}

Let $B_+$ and $B_-$ denote the corresponding products, pulled back to
$\mathbb S$.  Unimodular constants will never matter.

\section{The two-ended corona theorem}

The next estimate is the analytic core.  Notice that it is a product estimate,
not pairwise separation of zeros.

\begin{lemma}[Uniform one-sided product bounds]\label{lem:product-bound}
Under the hypotheses of Lemma~\ref{lem:bl-sum}, constants $c_+,c_->0$ exist
such that
\[
 |B_+(w)|\ge c_+\quad(|w|\ge1),\qquad
 |B_-(w)|\ge c_-\quad(|w|\le1).                       \tag{3.1}
\]
Here the products are written in the upper-half-plane coordinate.
\end{lemma}

\begin{proof}
For $a=a_\rho$ from the upper divisor, $|a|\le r<1$.  If $|w|=t\ge1$, then
\[
 |b_a(w)|\ge\frac{t-|a|}{t+|a|}\ge\frac{1-r}{1+r}=:q>0.  \tag{3.2}
\]
By (2.4), $\Im w\le t$, and $|w-\bar a|\ge t-r$, so
\[
       1-|b_a(w)|^2
       \le \frac{4t}{(t-r)^2}|a|
       \le C_r|a|,                                    \tag{3.3}
\]
where $C_r=\sup_{t\ge1}4t/(t-r)^2<\infty$.  On $[q,1]$ the ratio
$-\log x/(1-x^2)$ is bounded.  Therefore
\[
 -\log|B_+(w)|
 \le C_{q,r}\sum_{\rho\in\Sigma_+}m_\rho|a_\rho|,
\]
uniformly for $|w|\ge1$.  The right side is finite by Lemma~\ref{lem:bl-sum},
which proves the first bound.

The Möbius involution $w\mapsto-1/w$ preserves $\C_+$, interchanges the regions
$|w|\le1$ and $|w|\ge1$, and sends the lower divisor to a divisor satisfying the
upper estimates in (2.6).  Blaschke factors are preserved in modulus up to
unimodular constants.  The first argument therefore gives the second bound.
\end{proof}

\begin{theorem}[Two-ended corona]\label{thm:corona}
For a real-zero-free strip-admissible divisor,
\[
 \inf_{s\in\mathbb S}\bigl(|B_+(s)|+|B_-(s)|\bigr)
 \ge\min(c_+,c_-)>0.                                  \tag{3.4}
\]
Consequently there are $a,b\in H^\infty(\mathbb S)$ with
\[
                        aB_+ + bB_-=1.                \tag{3.5}
\]
\end{theorem}

\begin{proof}
Every $w\in\C_+$ lies in $|w|\ge1$ or $|w|\le1$.  Apply the corresponding
inequality in Lemma~\ref{lem:product-bound}.  Equation (3.5) is Carleson's
corona theorem, transported through $\Phi$.
\end{proof}

The argument allows arbitrary multiplicities and arbitrarily close pairs within
one half-plane.  The only separation is global and multiplicative: one product
is uniformly invertible on each of the two radial ends.

\subsection{A finite fixed sector}

Suppose $\Sigma_{\R}=\Sigma\cap\R$ is nonempty.  It is finite by (A1), and its
image under $\Phi$ lies on the compact semicircle $|w|=1$.  Let $B_{\R}$ be its
finite Blaschke product.

\begin{proposition}[Three-sector corona]\label{prop:three-corona}
The products $B_+,B_{\R},B_-$ are pairwise corona.  Hence the ideals they
generate in $H^\infty(\mathbb S)$ are pairwise comaximal.
\end{proposition}

\begin{proof}
The pair $(B_+,B_-)$ is covered by Theorem~\ref{thm:corona}.  On $|w|=1$ the
proof of Lemma~\ref{lem:product-bound} gives positive lower bounds for both
$B_+$ and $B_-$.  Thus each is nonzero at every zero of the finite product
$B_{\R}$.  Choose disjoint small neighbourhoods of those finitely many zeros.
On their union the relevant nonreal product is bounded below.  Off that union,
$B_{\R}$ is bounded below: on compact annuli this follows by continuity, while
at $0$ and infinity every factor of the finite product tends in modulus to one.
This proves the two remaining corona estimates.
\end{proof}

\section{Chinese remainders and quaternionification}

Assume from now on that the nonreal divisor is nonempty.  Put $B=B_+B_-$ and
define
\[
 \A_B=H^\infty(\mathbb S)/BH^\infty(\mathbb S),
 \qquad
 \Q_B=H^2(\mathbb S)/BH^2(\mathbb S).                \tag{4.1}
\]
The first is a commutative unital Banach algebra.  The second is a Hilbert
module over it.

\begin{proposition}[Topological Chinese remainder theorem]
\label{prop:crt}
Restriction to the two factors gives bounded isomorphisms
\[
 \begin{split}
 \A_B&\simeq
 H^\infty/B_+H^\infty\ \oplus\ H^\infty/B_-H^\infty,\\
 \Q_B&\simeq
 H^2/B_+H^2\ \oplus\ H^2/B_-H^2.
 \end{split}                                           \tag{4.2}
\]
If $a,b$ satisfy (3.5), then
\[
                   e=[bB_-]\in\A_B                 \tag{4.3}
\]
is the inverse image of $(1,0)$.  In particular $e^2=e$, and its class is
independent of the Bezout solution.
\end{proposition}

\begin{proof}
The Bezout identity implies comaximality.  The kernel of the restriction map is
the intersection of the two principal ideals.  If $f=B_+g=B_-h$, multiply
(3.5) by $f$ to obtain
\[
 f=aB_+B_-h+bB_-B_+g\in B_+B_-H^\infty.
\]
The reverse inclusion is immediate.  The same argument works in $H^2$ because
the Bezout coefficients are bounded multipliers.  It also constructs a bounded
inverse, proving (4.2).  Finally $bB_-=1-aB_+$, so (4.3) has residues $(1,0)$.
That residue pair determines the quotient class and proves both idempotency and
independence.
\end{proof}

Define two involutions on functions by
\[
       f^\#(s)=\overline{f(\bar s)},\qquad
       \sigma(f)(s)=f(1-s).                           \tag{4.4}
\]
The first is conjugate-linear and the second is linear.  Divisor invariance
makes both descend to $\A_B$ and $\Q_B$.

\begin{lemma}[Exchange law]\label{lem:exchange}
After harmless unimodular normalization of the two Blaschke factors,
\[
             e^\#=1-e,\qquad \sigma(e)=1-e.           \tag{4.5}
\]
\end{lemma}

\begin{proof}
Both involutions exchange the upper and lower divisors with multiplicity.  On
the Chinese remainder side, each therefore exchanges the two coordinates.
Since $e$ corresponds to $(1,0)$, its image corresponds to $(0,1)$, which is
$1-e$.
\end{proof}

Let the norm on $\Q_B$ be transported from the orthogonal direct sum in (4.2),
using the usual quotient Hilbert norms.  This is an equivalent Hilbert norm,
canonical after the ordered labels $+$ and $-$ have been fixed.  In this norm
$e$ is an orthogonal projection and $\#$ is antiunitary between the two
summands.

\begin{theorem}[Canonical quaternionic Hardy module]\label{thm:quaternion}
Set $m=2e-1$ and
\[
                         Jx=m x^\#.                   \tag{4.6}
\]
Then $m^2=1$, $m^\#=-m$, and $J$ is antiunitary for the Chinese remainder norm
with
\[
                         J^2=-1.                      \tag{4.7}
\]
For the original Hardy quotient norm, $J$ remains bounded and invertible.
\end{theorem}

\begin{proof}
Idempotency gives $(2e-1)^2=1$, while (4.5) gives
$m^\#=2(1-e)-1=-m$.  Hence
\[
 J^2x=m(m x^\#)^\#=mm^\#x=-m^2x=-x.
\]
In the direct sum (4.2), $m$ is the diagonal sign $(1,-1)$ and $\#$ is the
antiunitary exchange of summands.  Their product is antiunitary.  Equivalence of
the two quotient norms proves the last assertion.
\end{proof}

This is quaternionic in the precise operator-theoretic sense: $i$ from the
complex scalar structure and $J$ are anti-commuting real-linear isometries,
each squaring to $-1$.

\section{The oriented \texorpdfstring{$K$}{K}-class and its cyclic character}

For each $\rho\in\Sigma$, evaluation descends to a continuous character
$\ev_\rho:\A_B\to\C$.  Indeed $B(\rho)=0$, and the quotient norm controls the
value by taking infima over representatives.

\begin{theorem}[Infinite-order anti-Real class]\label{thm:kclass}
The virtual idempotent
\[
                 \kappa=[e]-[1-e]\in K_0(\A_B)       \tag{5.1}
\]
has infinite order.  Moreover
\[
                 \#_*(\kappa)=-\kappa,
                 \qquad \sigma_*(\kappa)=-\kappa.    \tag{5.2}
\]
\end{theorem}

\begin{proof}
If $\rho\in\Sigma_+$, then $\ev_\rho(e)=1$ and
$\ev_\rho(1-e)=0$.  The induced map
$K_0(\A_B)\to K_0(\C)=\Z$ sends $\kappa$ to $1$.  Therefore no nonzero integer
multiple of $\kappa$ vanishes.  Equations (5.2) follow from (4.5).
\end{proof}

Let $h:\Sigma\to\C$ satisfy
\[
                   \|h\|_{\ell^1(\Sigma,m)}
                   =\sum_{\rho\in\Sigma}m_\rho|h(\rho)|<\infty. \tag{5.3}
\]

\begin{proposition}[Zero cocycles]\label{prop:zero-cocycle}
The formula
\[
       \tau_h([f])=\sum_{\rho\in\Sigma}m_\rho h(\rho)f(\rho)    \tag{5.4}
\]
defines a continuous cyclic $0$-cocycle on $\A_B$, and
\[
 \langle[\tau_h],\kappa\rangle
 =\sum_{\rho\in\Sigma_+}m_\rho h(\rho)
  -\sum_{\rho\in\Sigma_-}m_\rho h(\rho).              \tag{5.5}
\]
\end{proposition}

\begin{proof}
For every representative $f+Bg$ of a quotient class the values at the divisor
are those of $f$.  Moreover $|f(\rho)|\le\|f\|_\infty$.  Taking the infimum over
representatives gives
\[
       |\tau_h([f])|\le \|h\|_{\ell^1(\Sigma,m)}\|[f]\|_{\A_B}.
\]
Thus (5.4) is well defined and continuous.  The algebra is commutative, so every
continuous linear functional is a cyclic $0$-cocycle.  Its matrix extension
pairs with an idempotent by evaluation of the ordinary matrix trace.  Since $e$
is one on $\Sigma_+$ and zero on $\Sigma_-$, while $1-e$ has the opposite
values, (5.5) follows.
\end{proof}

Mellin transforms of compactly supported smooth scaling tests give a natural
class of weights.  If
\[
                h_\varphi(s)=\int_\R\varphi(t)e^{ts}\dd t,
                \qquad \varphi\in C_c^\infty(\R),     \tag{5.6}
\]
then repeated integration by parts gives rapid decay in $|\Im s|$ on the closed
strip.  Polynomial zero counting implies (5.3).  Hence (5.5) is a distributional
oriented character of the scaling action.

\section{Real fixed points and the exact defect}

Let $B=B_+B_{\R}B_-$ and use Proposition~\ref{prop:three-corona}.  The
three-factor Chinese remainder theorem yields canonical idempotents
\[
 e_++e_{\R}+e_-=1,
 \qquad e_ie_j=\delta_{ij}e_i.                        \tag{6.1}
\]
Conjugation exchanges $e_+$ and $e_-$ and fixes $e_{\R}$.

\begin{theorem}[Fixed-sector defect law]\label{thm:defect}
Define
\[
        m_0=e_+-e_-,\qquad J_0x=m_0x^\#.             \tag{6.2}
\]
Then
\[
       m_0^2=1-e_{\R},\qquad m_0^\#=-m_0,
       \qquad J_0^2=-(1-e_{\R}).                      \tag{6.3}
\]
The class $\delta=[e_{\R}]\in K_0(\A_B)$ is fixed by $\#$, is nonzero if and
only if the real divisor is nonempty, and satisfies
\[
       \langle[\tau_h],\delta\rangle
       =\sum_{\rho\in\Sigma_{\R}}m_\rho h(\rho).      \tag{6.4}
\]
\end{theorem}

\begin{proof}
Orthogonality in (6.1) gives
\[
 (e_+-e_-)^2=e_++e_-=1-e_{\R}.
\]
The exchange law gives $m_0^\#=-m_0$, and the computation from
Theorem~\ref{thm:quaternion} yields (6.3).  If $\rho$ is a real zero, evaluation
at $\rho$ sends $[e_{\R}]$ to $1\in K_0(\C)$, so $\delta\ne0$.  The converse is
tautological because $e_{\R}=0$ when the factor is absent.  Formula (6.4) is
Proposition~\ref{prop:zero-cocycle} applied to the real idempotent.
\end{proof}

The next statement makes ``obstruction'' exact for the scalar cyclic module.

\begin{proposition}[Fixed-fibre obstruction]\label{prop:fibre}
Let $\rho\in\Sigma_{\R}$ and let
$\C_\rho$ be the one-dimensional $\A_B$-module defined by $\ev_\rho$.  There is
no conjugate-linear endomorphism $T$ of $\C_\rho$ satisfying $T^2=-1$.  Hence a
conjugate-linear, $\#$-semilinear operator squaring to $-1$ on the full cyclic
scalar quotient exists if and only if $e_{\R}=0$.
\end{proposition}

\begin{proof}
Every conjugate-linear endomorphism of a one-dimensional complex space has the
form $Tz=c\bar z$.  Its square is multiplication by $|c|^2$, which cannot equal
$-1$.  Any global $\#$-semilinear operator descends to the fixed evaluation
fibre, proving necessity.  If $e_{\R}=0$, Theorem~\ref{thm:quaternion} proves
sufficiency.
\end{proof}

Multiplicity does not evade the obstruction: evaluation at a real zero still
gives a one-dimensional quotient character, even when the local primary
component contains nilpotents.

\section{The completed zeta divisor}

Write
\[
 \xi(s)=\frac12s(s-1)\pi^{-s/2}\Gamma(s/2)\zeta(s).   \tag{7.1}
\]
Its nontrivial zero divisor in $\mathbb S$ satisfies all symmetries in (A3), and
the Riemann--von Mangoldt formula gives polynomial zero counting.  There are no
real zeros in $0<s<1$: the alternating eta series is positive there and
$\zeta(s)=\eta(s)/(1-2^{1-s})<0$.

It remains to identify the inner factor.  A function analytic across the finite
boundary of a strip can have singular inner mass only at the two prime ends.
Such mass has superexponential effect in the vertical coordinate.

\begin{lemma}[Endpoint singular factors]\label{lem:singular}
Let $F\in H^2(\mathbb S)$ extend holomorphically across every finite point of
$\partial\mathbb S$.  If its inner factor has nonzero singular mass at a prime
end, then on every fixed interior line $0<x<1$
\[
              \log|F(x+it)|\le -c e^{\pi|t|}+O(1)     \tag{7.2}
\]
in the corresponding direction, outside the zero set.
\end{lemma}

\begin{proof}
Transport $F$ to the upper half-plane by (2.2), and then to the disk.  Analytic
continuation across every finite boundary point forces the singular measure to
be supported at the images of $0$ and infinity.  A point mass $c>0$ contributes
the Poisson kernel of that endpoint to $\log|F|$.  The ray
$w=e^{-\pi t}e^{\pi ix}$ approaches $0$ non-tangentially for fixed $0<x<1$,
and the endpoint Poisson kernel there is comparable to $e^{\pi t}$; the other
prime end is identical after inversion.  Dividing an $H^2$ function by one of
its inner factors preserves its $H^2$ norm.  Point evaluation on a fixed
interior line has a bound independent of vertical translation.  The quotient
by the singular factor is therefore uniformly bounded on that line, which
gives (7.2).
\end{proof}

\begin{proposition}[The zeta inner factor is Blaschke]\label{prop:xi-inner}
The function $\xi$ belongs to $H^2(\mathbb S)\cap H^\infty(\mathbb S)$ and its
inner factor has no singular component.  Hence, up to a unimodular constant, it
is exactly $B=B_+B_-$.
\end{proposition}

\begin{proof}
Stirling's formula, standard polynomial bounds for $\zeta$ in vertical strips,
and the functional equation \cite[Chs.~IV and IX]{Titchmarsh1986} show that
$\xi$ is bounded and square-integrable
on the boundary lines, uniformly through the strip.  Thus it lies in both
Hardy classes.

If a singular endpoint mass were present, Lemma~\ref{lem:singular} would force
(7.2).  On the other hand, the Hardy--Littlewood mean-square theorem
\cite[Ch.~VII]{Titchmarsh1986} supplies a
sequence $t_n\to+\infty$ on which $|\zeta(1/2+it_n)|$ is bounded below by a
negative power of $t_n$; the same holds toward $-\infty$ by conjugation.
Stirling's formula then gives
\[
        \log|\xi(1/2+it_n)|\ge -C|t_n|\log(2+|t_n|),  \tag{7.3}
\]
which contradicts (7.2).  Therefore the singular factor is trivial.  The zeros
of the remaining inner factor are precisely the nontrivial zeros, with their
multiplicities, proving the assertion.
\end{proof}

\begin{corollary}[Unconditional Riemann quaternionification]
\label{cor:zeta}
For the complete nontrivial divisor of $\zeta$, the Hardy quotient
$H^2(\mathbb S)/BH^2(\mathbb S)$ carries the canonical quaternionic operator
$J$ of Theorem~\ref{thm:quaternion}; its oriented class $\kappa$ has infinite
order and all Mellin-test cyclic pairings are given by (5.5).  No assertion
about $\Re\rho$ is used.
\end{corollary}

\section{Mellin globalization of the scalar spectral quotient}

We now explain why the Hardy quotient is attached to the arithmetic source.
This section establishes a dense globalization, not an isomorphism of every
bornological model used in the literature.  We use the concrete
$\Hh_\cap,\Hh_-$ notation of \cite{Li2019}, which specializes Meyer's nuclear
Fr\'echet realization \cite{Meyer2005}.

Let
\[
       W=L^2\bigl(\R,(1+e^{2y})\dd y\bigr),           \tag{8.1}
\]
and for $g\in C_c^\infty(\R)$ define
\[
       (\cM g)(s)=\int_\R g(y)e^{sy}\dd y.            \tag{8.2}
\]
Plancherel on the two boundary lines gives
\[
 \|\cM g\|_{H^2_\partial}^2
 =2\pi\int_\R |g(y)|^2(1+e^{2y})\dd y.               \tag{8.3}
\]
Thus $\cM$ extends to a unitary map, up to $\sqrt{2\pi}$, from $W$ onto the
boundary-norm realization of $H^2(\mathbb S)$.  For completeness, surjectivity
is the strip Paley--Wiener theorem: inverse Fourier transformation of the
$\Re s=0$ boundary gives a function $g$, analyticity identifies the inverse
transform on $\Re s=1$ with $e^yg(y)$, and finiteness of the two boundary norms
is exactly $g\in W$.

Following \cite{Li2019}, set
\[
 \Hh_\cap=\{f\in\mathcal S(\R):f\text{ is even},\ f(0)=0,
                   \ \widehat f(0)=0\}.               \tag{8.4}
\]
Its compact logarithmic core carries the Fourier-moment condition
$\widehat f(0)=0$.  If $g(y)=f(e^y)$ on the positive half-line, this condition
is $\int e^yg(y)\dd y=0$.  Put
\[
       \D_0=\left\{g\in C_c^\infty(\R):
                    \int_\R e^yg(y)\dd y=0\right\}.   \tag{8.5}
\]
The other defining condition $f(0)=0$ is automatic on this compact logarithmic
core, whose even extension belongs to $\Hh_\cap$.  Each element of $\D_0$ is
also a strong Schwartz function on the multiplicative line.

\begin{lemma}[Moment-zero density]\label{lem:moment-density}
The subspace $\D_0$ is dense in $W$.  Consequently $\cM(\D_0)$ is dense in
$H^2(\mathbb S)$.
\end{lemma}

\begin{proof}
Choose $\psi\in C_c^\infty((0,1))$ with $\int\psi=1$, fix $R>0$, and set
$q_n(y)=e^{-y}n^{-1}\psi((y-R)/n)$.  Then
$\int e^yq_n(y)\dd y=1$, while
\[
 \|q_n\|_W^2
 \le n^{-1}(1+e^{-2R})\|\psi\|_2^2\longrightarrow0, \tag{8.6}
\]
because $q_n$ is supported in $(R,R+n)$.  Given $g\in C_c^\infty$, the
functions $g-(\int e^yg(y)\dd y)q_n$ lie in $\D_0$ and converge to $g$ in $W$.
Since $C_c^\infty$ is dense in $W$, the first assertion follows; (8.3) gives
the second.
\end{proof}

Set
\[
       G(s)=\frac12s(s-1)\pi^{-s/2}\Gamma(s/2),
       \qquad \xi=G\zeta.                             \tag{8.7}
\]
The function $G$ is a bounded zero-free multiplier in the open strip.  Its
inner factor has no Blaschke part.  Analytic continuation across the finite
boundary leaves only possible singular masses at the prime ends.  Such a mass
would force superexponential decay by Lemma~\ref{lem:singular}, whereas
Stirling gives, on every fixed interior line,
\[
             \log|G(x+it)|=-\frac{\pi}{4}|t|+O(\log(2+|t|)).
\]
Thus its inner factor is constant and $G$ is outer.

Let $\Hh_-$ be the strong Schwartz space on $\R_+^\times$; in logarithmic
coordinates it consists of smooth $g$ for which $e^{\alpha y}g(y)$ and all its
derivatives are Schwartz for every $\alpha\in\R$.  Define the source map on
this space by
\[
               U:\Hh_-\longrightarrow H^2(\mathbb S),
               \qquad Ug=G\cM g.                     \tag{8.8}
\]
Since $\D_0\subset\Hh_-$, Lemma~\ref{lem:moment-density} and outerness of $G$
show that $U(\Hh_-)$ is dense.  On $\Hh_\cap$ let $Z$ be the arithmetic
summation operator
\[
                    (Zg)(y)=\sum_{n\ge1}g(y+\log n).
\]
Poisson summation and the two vanishing conditions imply
$Z\Hh_\cap\subset\Hh_-$; this is the source-space mapping theorem used in the
Connes--Meyer realization \cite{Meyer2005,Li2019}.
On the common initial half-plane, followed by analytic continuation, it obeys
the Mellin multiplier law
\[
                        \cM(Zg)=\zeta\,\cM g.         \tag{8.9}
\]
This is the zeta operator in the Connes--Meyer quotient
$\Hh_-/Z\Hh_\cap$ after passage to logarithmic coordinates.

\begin{theorem}[Hardy globalization]\label{thm:globalization}
For $g\in\Hh_\cap$, the source map satisfies
\[
             U(Zg)=\xi\,\cM g,
\]
and
\[
       \overline{U(Z\Hh_\cap)}
       =\overline{\xi H^2(\mathbb S)}
       =BH^2(\mathbb S).                              \tag{8.10}
\]
It therefore induces a dense map
\[
                  \Hh_-/Z\Hh_\cap\longrightarrow
                  \Q_B=H^2/BH^2.
\]
This map intertwines logarithmic translations with the bounded multipliers
$e^{ts}$.
\end{theorem}

\begin{proof}
The first identity is (8.7)--(8.9).  Lemma~\ref{lem:moment-density} makes
$\cM(\D_0)$ dense in $H^2$.  The even lift of each element of $\D_0$ belongs
to $\Hh_\cap$, so this compact core gives the lower inclusion needed in (8.10).
Write the inner--outer factorization
$\xi=BO$.  Proposition~\ref{prop:xi-inner} identifies the inner factor as $B$,
and the outer factor $O$ has dense multiplier range.  Every Mellin transform of
an element of $\Hh_\cap$ lies in $H^2$, so the reverse inclusion is automatic.
Therefore
\[
 \overline{\xi\cM(\D_0)}=\overline{BOH^2}=BH^2
 =\overline{\xi\cM(\Hh_\cap)},
\]
which proves (8.10) and the dense quotient map.  Translation of $g$ by $t$ in
the logarithmic variable multiplies its Mellin transform by $e^{ts}$.  This
function is bounded on the strip for fixed $t$, so the action descends and the
intertwining follows.
\end{proof}

\begin{proposition}[Finite-packet character compatibility]
\label{prop:packet-character}
Let $F\subset\Sigma$ be a finite multiplicity-preserving packet invariant under
both divisor involutions, and let $B_F$ be its finite Blaschke product.  On
$H^2/B_FH^2$, the integrated scaling action against
$\varphi\in C_c^\infty(\R)$ has trace
\[
                       \sum_{\rho\in F}m_\rho h_\varphi(\rho). \tag{8.11}
\]
Its pairing with the packet restriction of $\kappa$ is the same sum with the
upper-minus-lower sign.
\end{proposition}

\begin{proof}
The finite Chinese remainder decomposition identifies the packet quotient with
the direct sum of local jet algebras
\[
             \bigoplus_{\rho\in F}\C[\varepsilon]/(\varepsilon^{m_\rho}).
\]
Multiplication by $e^{ts}$ is upper triangular on the monomial jet basis at
$\rho$, with constant diagonal entry $e^{t\rho}$.  Its trace on that block is
$m_\rho e^{t\rho}$.  Integration against $\varphi(t)$ gives (8.11).  The
corona idempotent is one on every upper block and zero on every lower block, so
pairing with $[e]-[1-e]$ inserts the stated signs.
\end{proof}

Letting invariant packets increase gives (5.5), because rapid decay of
$h_\varphi$ and polynomial zero counting give absolute convergence.  This is
the precise distributional compatibility with the Meyer/CCM spectral
character.  It makes no global Schatten assertion.

\section{Functional equation and cyclic transgression}

Let $\sigma$ be the automorphism induced by $s\mapsto1-s$ and form the algebraic
crossed product
\[
             \B_B=\A_B\rtimes_\sigma\Z/2.            \tag{9.1}
\]
It is generated by $\A_B$ and an element $u$ satisfying
\[
             u^2=1,\qquad uau=\sigma(a).              \tag{9.2}
\]
The same identities hold in the Banach crossed-product completion.

\begin{theorem}[Ordinary trace-blindness]\label{thm:traceblind}
Let $x=eu$ and $y=ue$.  Then
\[
        xy=e,qquad yx=1-e,                            \tag{9.3}
\]
so the image of $\kappa$ in $K_0(\B_B)$ is zero.  Moreover
\[
        m=2e-1=xy-yx.                                  \tag{9.4}
\]
Consequently every ordinary cyclic $0$-cocycle $\tau$ on $\B_B$ satisfies
$\tau(m)=0$, and $m$ represents zero in ordinary $HC_0(\B_B)$.
\end{theorem}

\begin{proof}
Using $e^2=e$, $u^2=1$, and $ueu=\sigma(e)=1-e$, we obtain
\[
 xy=(eu)(ue)=e^2=e,
 \qquad
 yx=(ue)(eu)=ueu=1-e.
\]
Thus $x$ and $y$ implement algebraic equivalence of the two idempotents, which
proves the $K_0$ assertion.  Subtracting the two identities gives (9.4).  A
cyclic $0$-cocycle is a trace, so it kills commutators.  The last assertion is
the defining quotient description of $HC_0$.
\end{proof}

The vanishing is exact but not total.  Equation (4.5) also gives
\[
                       mu=-um.                        \tag{9.5}
\]
Indeed $eu=u(1-e)$, and direct expansion proves (9.5).  Since $m^2=u^2=1$,
the two elements satisfy the split rank-two Clifford relations.  The inner
automorphism $\gamma=\operatorname{Ad}(m)$ grades $\B_B$: the coefficient
algebra is even and $u$ is odd.

Fix $\rho\in\Sigma_+$ and write $\rho'=1-\rho\in\Sigma_-$.  There is a
two-dimensional representation
\[
 \pi_\rho(a)=
 \begin{pmatrix}a(\rho)&0\\0&a(\rho')\end{pmatrix},
 \qquad
 \pi_\rho(u)=
 \begin{pmatrix}0&1\\1&0\end{pmatrix}.              \tag{9.6}
\]
Then
\[
 \pi_\rho(e)=
 \begin{pmatrix}1&0\\0&0\end{pmatrix},
 \qquad
 \pi_\rho(m)=
 \begin{pmatrix}1&0\\0&-1\end{pmatrix}.             \tag{9.7}
\]

\begin{theorem}[Graded cyclic survivor]\label{thm:graded}
Define
\[
          \Str_\rho(T)=\Tr\bigl(\pi_\rho(m)T\bigr).  \tag{9.8}
\]
On the graded image of $\B_B$, $\Str_\rho$ is a graded cyclic $0$-cocycle and
\[
                       \Str_\rho(\pi_\rho(m))=2.      \tag{9.9}
\]
Therefore $m$ is nonzero in graded $HC_0(\B_B)$ even though it is zero in
ordinary $HC_0(\B_B)$.
\end{theorem}

\begin{proof}
The matrix $\pi_\rho(m)$ is the grading operator in (9.6)--(9.7).  The standard
supertrace identity gives
\[
 \Str_\rho(ST)=(-1)^{|S||T|}\Str_\rho(TS)
\]
for homogeneous matrices, so (9.8) is a graded cyclic cocycle.  Finally
\[
 \Str_\rho(\pi_\rho(m))
 =\Tr(\pi_\rho(m)^2)=\Tr(I_2)=2.
\]
A graded boundary pairs trivially with every graded cocycle.  The nonzero
pairing proves the last assertion.
\end{proof}

\begin{definition}
For the exchanged idempotent $e$, we call the passage
\[
 [e]-[1-e]\in K_0(\A_B)
 \quad\longmapsto\quad
 [2e-1]_{\mathrm{gr}}\in HC^{\mathrm{gr}}_0(\A_B\rtimes_\sigma\Z/2)
                                                               \tag{9.10}
\]
the \emph{Real-corona cyclic transgression}.  This terminology names the
explicit class construction proved above; it does not assert a general
spectral-sequence differential for arbitrary crossed products.
\end{definition}

Combining Theorems~\ref{thm:traceblind} and \ref{thm:graded} yields the promised
information law:
\[
\begin{array}{ccl}
\text{anti-Real orientation in }K_0(\A_B)
 &\xrightarrow{\ \rtimes\sigma\ }&0\text{ in ordinary }K_0,\\[1mm]
m\text{ in ordinary }HC_0(\B_B)
 &=&0,\\[1mm]
m\text{ in graded }HC_0(\B_B)
 &\ne&0.
\end{array}                                             \tag{9.11}
\]
The functional equation does not merely impose a symmetry on the character;
when internalized, it changes the homology theory in which orientation is
visible.

\section{Formal verification of the algebraic core}

The algebraic implications used in Sections 4, 6, and 9 have been formalized
unconditionally in Lean~4 with mathlib.  The source file is
\begin{center}
\texttt{formal/GaloisForLFunctions/RHTower/RealCoronaQuaternionicCyclic.lean}.
\end{center}
It proves, for arbitrary rings and compatible Real modules where possible:
\begin{enumerate}
\item $(2e-1)^2=1$ and $\sigma(2e-1)=-(2e-1)$;
\item $(m\sigma)^2=-1$ and construction of the odd Real structure;
\item $(e_+-e_-)^2=1-e_{\R}$ and the partial square law (6.3);
\item the equivalences (9.3), commutator identity (9.4), and annihilation by
      every additive cyclic trace;
\item anticommutation (9.5) and the explicit matrix supertrace value $2$.
\end{enumerate}
The file contains no \texttt{sorry}, \texttt{admit}, added axiom,
\texttt{opaque}, or \texttt{unsafe} declaration.  It passes direct elaboration,
the repository formal lint gate, and the targeted module build.  Analytic
existence of the corona idempotent is proved in Sections 2--3 and is not encoded
as an axiom in the formalization.

\section{Relation to prior constructions}

The adeles class-space cohomology of Connes--Consani--Marcolli is a cokernel in
cyclic modules, designed so that the scaling action has Weil's explicit formula
as a Lefschetz character \cite{CCM2007}.  Our construction does not replace that
cohomology.  Theorem~\ref{thm:globalization} supplies a Hilbert Hardy completion
of its scalar zeta multiplier quotient and equips that completion with a
bounded idempotent unavailable at the purely distributional level.

Meyer's representation is nuclear in a bornological sense and realizes the
full Weil distribution \cite{Meyer2005}.  The present cyclic cocycles agree with
its zero character on finite packets and by absolutely convergent test-function
limits.  This is exactly the compatibility proved here.  A Hilbert trace-class
upgrade would require separate operator-ideal estimates and is neither needed
nor claimed.

Paugam's spectral pairing encodes the functional equation by a symplectic or
orthogonal duality, depending on the $L$-function \cite{Paugam2009}.  The new
ingredient here is the corona idempotent: it makes the upper/lower sign a
bounded canonical multiplier, produces an explicit $K_0$ class, and identifies
the real fixed sector as an obstruction rather than deleting it.  The crossed
product then shows why the same orientation is invisible to ordinary traces
once functional symmetry becomes inner.

General cyclic homology of crossed products is governed by mixed complexes and
conjugacy-sector decompositions \cite{CarboniGuccione2009}.  Theorem
\ref{thm:graded} is more specific: it identifies the concrete arithmetic
orientation element, proves its ordinary commutator vanishing, and exhibits the
minimal graded cocycle that still detects it.

Connes and Consani's $\zeta$-cycles use Hochschild homology and trace maps to
organize stable spectral realizations \cite{ConnesConsani2023}.  The class
$\kappa$ and its transgression give a complementary Real/cyclic invariant on
the Hardy globalization.  They depend only on the complete divisor and its two
canonical involutions, not on a finite prolate truncation.

\section{Consequences and open problems}

The theorem package separates three questions often conflated in spectral
approaches to zeta.

\begin{enumerate}
\item \emph{Can the upper and lower spectral halves be separated by a bounded
multiplier?}  Yes, unconditionally, by the two-ended corona theorem.
\item \emph{Does the resulting Real module have a quaternionic sign?}  Yes
exactly off the fixed real sector; the defect is the fixed idempotent
$e_{\R}$.
\item \emph{Does internalizing the functional equation preserve the ordinary
cyclic orientation?}  No.  It kills the ordinary $K_0$ and $HC_0$ classes, but
the orientation survives canonically in graded cyclic degree zero.
\end{enumerate}

Several hard problems remain.

\begin{enumerate}[label=\textup{(P\arabic*)}]
\item Determine whether the graded class $[m]$ admits a local index formula on
      an unbounded zeta spectral triple compatible with the complete divisor.
\item Lift $\kappa$ from the scalar Hardy globalization to the full adeles
      class-space cyclic module and compare its bivariant Chern character with
      the CCM Lefschetz correspondence.
\item Classify completed automorphic $L$-functions for which the Real-corona
      idempotent is smooth for a natural dense holomorphic subalgebra of the
      crossed product.  Smoothness, unlike boundedness, is the missing input
      for higher entire cyclic pairings.
\item Compute the fixed-sector defect for self-dual families in which central
      real zeros vary, and compare $[e_{\R}]$ with analytic rank.
\end{enumerate}

These problems are independent of RH.  They concern whether the bounded
orientation constructed here participates in deeper noncommutative-geometric
index theory.

\begin{thebibliography}{99}

\bibitem{CarboniGuccione2009}
G. Carboni, J. A. Guccione, and J. J. Guccione,
\emph{Cyclic homology of crossed products},
J. Algebra \textbf{321} (2009), 2827--2860;
\href{https://arxiv.org/abs/0805.0582}{arXiv:0805.0582}.

\bibitem{Connes1994}
A. Connes,
\emph{Noncommutative Geometry},
Academic Press, San Diego, 1994.

\bibitem{Connes1999}
A. Connes,
\emph{Trace formula in noncommutative geometry and the zeros of the Riemann
zeta function},
Selecta Math. (N.S.) \textbf{5} (1999), 29--106.

\bibitem{CCM2007}
A. Connes, C. Consani, and M. Marcolli,
\emph{The Weil proof and the geometry of the adeles class space},
in \emph{Noncommutative Geometry, Arithmetic, and Related Topics},
Johns Hopkins Univ. Press, 2011, 339--405;
\href{https://arxiv.org/abs/math/0703392}{arXiv:math/0703392}.

\bibitem{ConnesConsani2023}
A. Connes and C. Consani,
\emph{Hochschild homology, trace map and $\zeta$-cycles},
Chinese Phys. C \textbf{47} (2023), 013106;
\href{https://arxiv.org/abs/2207.10419}{arXiv:2207.10419}.

\bibitem{CCM2025}
A. Connes, C. Consani, and H. Moscovici,
\emph{Zeta spectral triples}, preprint (2025),
\href{https://arxiv.org/abs/2511.22755}{arXiv:2511.22755}.

\bibitem{Garnett2007}
J. B. Garnett,
\emph{Bounded Analytic Functions}, revised first edition,
Graduate Texts in Mathematics 236, Springer, New York, 2007.

\bibitem{LopattoRochberg2016}
P. Lopatto and R. Rochberg,
\emph{Schatten-class truncated Toeplitz operators},
Proc. Amer. Math. Soc. \textbf{144} (2016), 637--649.

\bibitem{Li2019}
X.-J. Li,
\emph{On spectral theory of the Riemann zeta function},
Sci. China Math. \textbf{62} (2019), 2317--2330.

\bibitem{Meyer2005}
R. Meyer,
\emph{On a representation of the idele class group related to primes and zeros
of $L$-functions},
Duke Math. J. \textbf{127} (2005), 519--595.

\bibitem{Nikolski2002}
N. K. Nikolski,
\emph{Operators, Functions, and Systems: An Easy Reading. Vol. 2},
Mathematical Surveys and Monographs 93, American Mathematical Society,
Providence, 2002.

\bibitem{Paugam2009}
F. Paugam,
\emph{Sym\'etries spectrales des fonctions z\^etas},
J. Th\'eor. Nombres Bordeaux \textbf{21} (2009), 713--720;
\href{https://arxiv.org/abs/0803.0199}{arXiv:0803.0199}.

\bibitem{Titchmarsh1986}
E. C. Titchmarsh,
\emph{The Theory of the Riemann Zeta-Function}, second edition, revised by
D. R. Heath-Brown, Oxford University Press, Oxford, 1986.

\end{thebibliography}

\end{document}
